% Source-led account of the implemented Chen 1+2 route; local notation only.
\section{Chen's theorem: from a signed sieve weight to representations}
\label{bp:chen-route}

A prime $p<N$ determines its partner $N-p$ uniquely.  We count representations
by the prime $p$, rather than by factorizations of its partner.  The public finite set is
\[
 G(N)=\{p<N:p\text{ prime},\ 2\le N-p,\ \Omega(N-p)\le2\},
\]
Thus $|G(N)|=R_2(N)$ in the overview notation; $\Omega(n)$ counts prime factors with multiplicity.
Thus a prime partner and a square of a prime are both permitted, while
$N-p=1$ is excluded.  This is exactly \texttt{chenGoodRepresentations}.
The quantitative target is $|G(N)|\ge0.67 \mathcal X_N$ for all sufficiently large
even $N$, with the normalization
\[
 \mathcal X_N=\mathfrak S_{\mathrm{Liu}}(N)\frac{N}{(\log N)^2},\qquad
 \mathfrak S_{\mathrm{Liu}}(N)=\prod_{\substack{r\mid N\\r>2\ \text{prime}}}\frac{r-1}{r-2}
                 \prod_{r>2\ {\rm prime}}\left(1-\frac1{(r-1)^2}\right).
\]
The second product is the positive constant $U=C_2$ from the foundations chapter, and $\mathfrak S_{\mathrm{Liu}}(N)\ge U>0$.  This uniform lower bound will let us pay errors
without choosing a different normalization for each even integer.

\subsection{A finite weight that recognizes the desired partners}
Put $z=\max(2,\lfloor N^{1/10}\rfloor)$ and $y=\lceil N^{1/3}\rceil$.
The corrected candidate set is
\[
 \mathcal C(N)=\{p<N:p\text{ prime},\ N-p\ge2,
          \ r\nmid N-p\text{ for every prime }r<z\}.
\]
Sifting has removed small factors, but a surviving partner can still have
three or more prime factors. Let $\mathcal Q_{z,y}$ be the primes in $[z,y)$.
For a prime $r$, let $\nu_r(n)$ be its exponent in the factorization of $n$. To charge these survivors, define
\begin{align*}
 v(n)&=\sum_{r\in\mathcal Q_{z,y},\,r\mid n}\nu_r(n),\\
 t(n)&=\#\{r:z\le r<y,\ r\text{ prime},\
       \exists s,u\text{ prime},\ y\le s\le u,\ r<s,\ rsu=n\},\\
 w(n)&=1-\tfrac12v(n)-\tfrac12t(n).
\end{align*}
The second count records first factors $r$, with the larger factors existentially quantified.  In particular, its definition includes $s=u$.
For a $z$-rough integer $1\le n<y^3$, meaning that no prime below $z$ divides $n$, positivity of $w(n)$ forces
$n=1$, a prime, or a product of two primes.  Indeed, the nonnegative
integer $v(n)+t(n)$ must be at most one.  If $v(n)=0$, all factors
are at least $y$, so three factors contradict $n<y^3$.  If $v(n)=1$,
there is exactly one medium factor, with valuation one; two remaining
large factors would contribute to $t(n)$ and use up the second unit.
\inputleannode{bp:chen-positive-weight}
For $N\ge9$ we have $z<y$ and $N\le y^3$; a candidate's positive
prime component gives the strict inequality $N-p<y^3$.
Consequently a bad candidate has penalty $v(N-p)+t(N-p)\ge2$.
A good candidate has $w(N-p)\le1$, so summing the pointwise comparison
$w(N-p)\le\mathbf1_{G(N)}(p)$ gives
\[\begin{gathered}
 |G(N)|\ge|\mathcal C(N)|-\tfrac12(P(N)+T(N))=W(N)-\tfrac12T(N),\\
 P(N)=\sum_{p\in\mathcal C(N)}v(N-p),\qquad
 T(N)=\sum_{p\in\mathcal C(N)}t(N-p),\\
 W(N)=|\mathcal C(N)|-\tfrac12P(N).
\end{gathered}\]
Here $W$ is \texttt{jurkatRichertWeightedCount}; $P,T$ are the corrected prime-power and triple penalties.  The factor $1/2$ comes
from the bad-candidate inequality, not from a symmetry of a switched sum.
\inputleannode{bp:chen-finite-counting}

\subsection{The source weight separates a lower sieve from upper sieves}
The analytic lower bound first treats a slightly different finite carrier,
\[
 \mathcal S(N)=\{p\le N:p\text{ prime},\
    r\nmid N-p\text{ for every odd prime }r\le N^{1/10}\}.
\]
Let $\mathcal Q(N)$ be the primes $N^{1/10}<q\le N^{1/3}$, and set
\[
 H(N)=\sum_{q\in\mathcal Q(N)}\#\{p\in\mathcal S(N):q\mid N-p\},
 \qquad W_{\rm src}(N)=|\mathcal S(N)|-\tfrac12 H(N).
\]
Finite interchange of sums identifies this with the source weight
$\sum_{p\in\mathcal S(N)}(1-\frac12\#\{q\in\mathcal Q(N):q\mid N-p\})$.
Each medium prime is charged once, irrespective of valuation.  A lower bound for $|\mathcal S(N)|$ and an upper bound for $H(N)$ therefore
advance the same signed inequality in the required directions.

For the base sieve, remove the exceptional unsifted primes whose complement
is divisible by an odd cutoff prime dividing $N$.  Such an exceptional
prime equals that divisor.  The remaining sifting product and main mass are
\[\begin{gathered}
 \mathcal P_N=\prod_{\substack{2<r\le N^{1/10}\\r\ \text{prime},\ r\nmid N}}r,
 \qquad X=\operatorname{li}_{2/\log2}(N),\\
 V_N=\prod_{\substack{r\mid\mathcal P_N\\r\ \text{prime}}}\left(1-\frac1{r-1}\right),
\end{gathered}\]
Here $\operatorname{li}_{\kappa}=L_{\kappa}$ in the foundations notation:
\[
 \operatorname{li}_{\kappa}(x)=\kappa+\int_2^x\frac{dt}{\log t}\quad(x>1),
 \qquad \operatorname{li}_{\kappa}(x)=\kappa\quad(0\le x\le1).
\]
The latter is the implementation's totalized-integral convention, not a
Cauchy principal value. In particular $X=L_*(N)$, where
$L_*=L_{2/\log2}$ and $L_*(0)=L_*(1)=2/\log2$.
See \sourcefile{MathlibNt/SieveTheory/Arithmetic/LiuLogarithmicIntegral.lean}{genuine logarithmic integral}
and \sourcefile{MathlibNt/SieveTheory/Distribution/LiuPan/LiuPanCombinedAbelDeterministic.lean}{finite-endpoint convention}.
The density is $1/(r-1)$; $R_N(d)=|\mathcal A_d|-X/\varphi(d)$ is the discrepancy of the unsifted complement
carrier: $\mathcal A_d$ consists of the remaining complements divisible by $d\mid\mathcal P_N$.
With level $D=N^{1/2-\epsilon}$, lower Rosser coefficients give a finite
lower bound by their density sum minus
$\mathcal R_0=\sum_{d\mid\mathcal P_N,\ d\le D}|R_N(d)|$.
For the explicit sieve functions put
\[
 I(u)=\int_2^{u-1}\frac{\log(t-1)}t\,dt,\qquad
 J=\int_3^4\frac{I(u)}u\,du,\qquad
 K=\int_3^4\frac{10I(u)}{u(5-u)}\,du.
\]
The implemented density theorem yields, for each $\delta>0$, a fixed
$\epsilon>0$ and eventually for even $N$
\[
 |\mathcal S(N)|\ge (f(5)-\delta)XV_N-\mathcal R_0,
 \qquad f(5)=\frac{2e^\gamma}{5}(\log4+J).
\]

\inputleannode{bp:chen-lower-sieve}
The density input is supplied by the constructed Jurkat--Richert delay
functions through the modern Suzuki varying-family comparison.  Continuity
at $s=5$ permits a fixed small loss $s=5-10\epsilon$.
Ordinary Bombieri--Vinogradov pays $\mathcal R_0$ at this power-reduced
level; the removed exceptional primes receive a separate power-saving bound.
Mertens normalization supplies the lower estimate
$XV_N\ge(20e^{-\gamma}-\delta)\mathcal X_N$ for any fixed positive margin.
Combining the three eventual bounds gives, for every $\eta>0$,
\[
 |\mathcal S(N)|\ge (B_0-\eta)\mathcal X_N,\qquad B_0=8(\log4+J).
\]
The proof chooses the density loss first, then invokes distribution at that same $\epsilon$; all thresholds are intersected before choosing $N$.
\inputleannode{bp:chen-base-asymptotic}

\subsection{The varying-level upper penalty and its distribution bill}
For each $q\in\mathcal Q(N)$, condition the same complement carrier by
$q\mid N-p$.  Its nominal mass is $X/(q-1)$, including the lane $q\mid N$;
the exact discrepancy, rather than a reduced-residue approximation, handles
that lane.  The upper sieve uses
\[
 D_q=N^{1/2-\epsilon}/q,\qquad
 s_q=\frac{\log D_q}{\log N^{1/10}}
     =5-10\epsilon-10\frac{\log q}{\log N}.
\]
Integer support is $d<\lfloor D_q\rfloor+1$.  Summing the finite upper
Rosser inequalities gives $H\le\mathcal M_q+\mathcal R_q$, with
$\mathcal M_q$ the sum of $X/(q-1)$ times the actual upper density sums,
and $\mathcal R_q$ the sum of their absolute-coefficient remainder bills.
\inputleannode{bp:chen-finite-upper}
The modern all-depth Suzuki theorem controls these density sums uniformly
in the conditioned sieve.  In the range used here its continuous factor is
\[
 F(s)=\begin{cases}2e^\gamma/s,&s\le3,\\
 (2e^\gamma/s)(1+I(s)),&s>3,\end{cases}

\]
Thus the density model is $XV_N\sum_{q\in\mathcal Q(N)}F(s_q)/(q-1)$.
Partial summation retains both $s_q$ and the exact $q-1$.  For any margin, one chooses $0<\epsilon<1/60$ and then
obtains the coefficient $B_1=8(\log8+K/2)$, using the integrals defined above.
Distribution is needed for the combined modulus $qd$.  On $q\nmid N$
the finite bill is dominated by
\[
 \sum_{\substack{q\in\mathcal Q(N)\\q\nmid N}}
 \sum_{\substack{d\mid\mathcal P_N\\d\le D_q}}
       3^{\omega(d)}E_{\rm prime}(N,qd),
\]
where the fixed-endpoint error and the foundations' prefix error are
\[\begin{gathered}
 E_{\rm prime}(N,m)=\max_{l\in(\mathbb Z/m\mathbb Z)^\times}
   \left|\pi(N;m,l)-\frac{L_*(N)}{\varphi(m)}\right|,\\
 E^*(N,m)=\max_{n\in\{0,\ldots,N\}}E_{\rm prime}(n,m),
 \qquad E_{\rm prime}(N,m)\le E^*(N,m).
\end{gathered}\]
Here $\pi(N;m,l)$ counts primes at most $N$ in residue class $l$ modulo $m$.
Modulo one the canonical residue is zero; for modulus zero both maxima
are defined to be zero. Thus this is the same genuine-$L_*$ normalization
as ordinary BV, but only the endpoint maximum is consumed in the displayed bill.
The definitions and endpoint-to-prefix comparison are in
\sourcefile{MathlibNt/SieveTheory/Distribution/BombieriVinogradov.lean}{ordinary AP errors};
the exact weighted bill is in
\sourcefile{MathlibNt/SieveTheory/Switching/VaryingPrimeSieve.lean}{combined-modulus error sum}.
The map $(q,d)\mapsto qd$ is injective here: $q$ is above every prime
factor of $d$.  Finite Cauchy--Schwarz pays the weight with a
$9^{\omega}/d$ moment; ordinary BV is requested with saving $2A+10$
to leave saving $A$.  Nonreduced lanes and exceptional support corrections
are bounded separately and then absorbed into $\mathcal X_N$.
\inputleannode{bp:chen-weighted-bv}
The upper asymptotic allocates one third of its margin to density comparison,
one third to prime partial summation, and one third to distribution.
It concludes $H(N)\le(B_1+\eta)\mathcal X_N$ eventually for even $N$.
\inputleannode{bp:chen-upper-asymptotic}
Using a base margin $\eta/2$ and an upper margin $\eta$ now gives
\[\begin{gathered}
 W_{\rm src}(N)\ge(2.6408-\eta)\mathcal X_N,\\
 B_0-\tfrac12B_1=8\bigl(\log4-\tfrac12\log8+J-K/4\bigr)\ge2.6408.
\end{gathered}\]
The last comparison uses the proved $J-K/4\ge-0.0164725$ and a certified
lower bound for $\log2$; it is an integral inequality, not floating-point quadrature.
\inputleannode{bp:chen-distinct-lower}

\subsection{Paying for rounded endpoints and repeated prime factors}
Write $W_{\rm dist}=|\mathcal C|-\frac12\sum_{p\in\mathcal C}
\#\{r\in[z,y):r\text{ prime},\ r\mid N-p\}$.
For even $N>2^{110}$, the implemented source comparison is
$W_{\rm src}-10N^{9/10}\le W_{\rm dist}$.
It splits the two carriers into their intersection and exceptional fibres:
the lower-floor fibre, $p=2$, and the unit complement $N-p=1$.
Signed weights need not be monotone under enlarging the candidate set, so this comparison is necessary.
Repeated medium factors have their own finite estimate:
\[
 P(N)\le\sum_{p\in\mathcal C(N)}\#\{r\in[z,y):r\text{ prime},\ r\mid N-p\}
           +60N^{9/10}.
\]
Divisibility by $r^2$ is bounded by counting quotients, giving at most
$6N^{9/10}$ candidate--prime pairs; each relevant valuation is at most ten.
\inputleannode{bp:chen-prime-powers}
Hence $W_{\rm dist}-30N^{9/10}\le W$.
Since $60N^{9/10}\le\rho \mathcal X_N$ eventually for every $\rho>0$,
the source-to-corrected loss and the valuation loss fit within an arbitrary
coefficient margin.  Taking the source margin and the power-error margin
to be $\eta/2$ proves $W(N)\ge(2.6408-\eta)\mathcal X_N$.
This lower bound has already paid the entire prime-power penalty $P$.
\inputleannode{bp:chen-valuation-lower}

\subsection{Moving the remaining triple penalty into a Selberg square}
Only $T$ remains to be subtracted.  Put $z_0=\lfloor N^{1/10}\rfloor$,
$y_0=\lfloor N^{1/3}\rfloor$, and let $\mathcal L$ consist of prime pairs
$(r,s)$ satisfying $z_0<r\le y_0<s$ and $rs^2\le N$.
Liu's characteristic weight $b(a)$ is one exactly when $a=rs$ for such
a pair; it is unique, and sums involving $b$ run over its support.  On $z\le r<y\le s$,
let $T_{\rm src}$ count $(r,s,u)$ with $(r,s)\in\mathcal L$,
$u$ prime, $s\le u$, $rsu\le N$, and $N-rsu\in\mathcal C(N)$.
A first-factor witness counted by $T$ chooses one ordered pair $(s,u)$.
The map to $(r,s,u)$ recovers the candidate as $p=N-rsu$, so it is injective.
Its only cutoff exceptions are $r=z_0$ and $s=y_0$; quotient counting pays
both, giving $T\le T_{\rm src}+13N^{9/10}$ once $z_0\ge2$.
\inputleannode{bp:chen-triple-endpoints}
For even $N$ and $R=\lfloor N^{1/4-\epsilon/2}\rfloor\ge1$, put
$Q=\prod_{r\le R,\ r\ \text{prime},\ r\nmid N}r$ and
$\mathcal D=\{d\mid Q:d\le R\}$.  The Selberg square is
\[
 \mathcal B_\lambda(N)=\sum_a b(a)
   \sum_{\substack{u\ \text{prime}\\au\le N}}
       \left(\sum_{\substack{d\in\mathcal D\\d\mid N-au}}\lambda_d\right)^2.
\]
The optimal coefficients satisfy $\lambda_1=1$, $|\lambda_d|\le1$, and vanish off $\mathcal D$.  If the candidate prime $p=N-rsu$ does not
divide $Q$, its divisor packet contains only $d=1$ and its square is one.
Primes $p\mid Q$ contribute to a residual $E_Q$.  Nonnegativity allows us to drop $s\le u$ and enlarge the common pair region,
yielding $T_{\rm src}\le\mathcal B_\lambda+E_Q$.
\inputleannode{bp:chen-square-majorant}
For each source pair the exceptional residual injects into the prime divisors
of $Q$.  Pair support has size at most $\lceil N^{2/3}\rceil+1$, whereas
$Q$ has at most $R+1$ prime divisors.  Thus $E_Q\le6N^{11/12}$ for
$N\ge1$ and $\epsilon\ge0$.  Both this loss and the endpoint loss are
smaller than $C_A N/(\log N)^A$ for every fixed $A>0$.

\subsection{The square's main term and the switched distribution error}
Expanding the square replaces simultaneous divisibility by
$[d_1,d_2]\mid N-au$.  Finite interchange gives exactly
\[\begin{gathered}
 \mathcal B_\lambda=M_1+\mathcal E_\lambda,\\
 M_1=\left(\sum_{d_1,d_2\in\mathcal D}
     \frac{\lambda_{d_1}\lambda_{d_2}}{\varphi([d_1,d_2])}\right)
       \sum_a b(a)\operatorname{li}_2(N/a).
\end{gathered}\]
The error uses the same coefficients and the actual progression discrepancy
$\pi(N;a,d,l)-\operatorname{li}_2(N/a)/\varphi(d)$, where
$\pi(N;a,d,l)$ counts primes $u$ with $au\le N$ and $au\equiv l\pmod d$.
\inputleannode{bp:chen-square-expansion}
The optimizer uses reverse-divisor M\"obius inversion of the diagonal coordinates $\mu(\ell)g(\ell)/\mathcal G$, where
$g(\ell)=\prod_{\substack{r\mid\ell\\r\ \text{prime}}}(r-2)^{-1}$ and
$\mathcal G=\sum_{\ell\in\mathcal D}g(\ell)$.
Diagonalizing the quadratic form proves its value is exactly $1/\mathcal G$.
The denominator estimate bounds this by
$(8+\delta)\mathfrak S_{\mathrm{Liu}}(N)/\log N$ whenever
$\delta>0$ and $0<\epsilon<\delta/(2(8+\delta))$, eventually along even $N$.
The other factor must be estimated with the genuine logarithmic integral,
not identified with its proxy $x/\log x$.

Write the exact prime-pair reciprocal-log mass as
\[
 S_{\rm rec}(N)=\sum_{(r,s)\in\mathcal L}
       \frac{1}{rs\log(N/(rs))}.
\]
In logarithmic coordinates $\alpha=\log r/\log N$ and
$\beta=\log s/\log N$, the limiting domain is
\[
 \mathcal D_{\rm Liu}=\{(\alpha,\beta):1/10\le\alpha\le1/3,
                  \ 1/3\le\beta\le(1-\alpha)/2\}.
\]
The sloping edge is precisely the condition $rs^2\le N$;
the finite sum still uses the rounded cutoffs defining $\mathcal L$.
The limiting prime reciprocal measures give the factors $d\alpha/\alpha$
and $d\beta/\beta$, while
$\log(N/(rs))=(1-\alpha-\beta)\log N$. The source integral is therefore
\[\begin{aligned}
 I_{\rm Liu}
 &=\int_{1/10}^{1/3}\frac{d\alpha}{\alpha}
       \int_{1/3}^{(1-\alpha)/2}\frac{d\beta}{\beta(1-\alpha-\beta)}\\
 &=\int_{1/10}^{1/3}
       \frac{\log(2-3\alpha)}{\alpha(1-\alpha)}\,d\alpha
 <0.49254.
\end{aligned}\]
All denominators are positive on this domain. Partial fractions evaluate
the inner integral as $\log(2-3\alpha)/(1-\alpha)$, proving the reduction.
The strict constant bound is proved for this integral itself: after the
change $t=(1-3\alpha)/(3-3\alpha)$ it becomes
\[
 3\int_0^{7/27}\frac{\log((1+t)/(1-t))}{1-3t}\,dt.
\]
On this interval the source bounds the logarithm above by
\[
 P_{\log}(t)=2\sum_{j=0}^{5}\frac{t^{2j+1}}{2j+1}
                    +\frac{729}{340}t^{13}.
\]
This follows from the logarithmic series with its bounded positive tail.
Exact integration of $3P_{\log}(t)/(1-3t)$, followed by a rational upper
bound for $\log(9/2)$, gives the strict inequality above. This is not a
comparison of two preassigned decimals or a floating-point quadrature;
see \sourcefile{MathlibNt/SieveTheory/Liu/Weights/LiuWeightMainIntegral.lean}{source integral and strict bound}.

For each fixed $\tau>0$, the actual transfer proves
\[
 \log N\,S_{\rm rec}(N)\le I_{\rm Liu}+\tau
 \quad\hbox{for all sufficiently large }N.
\]
It first fixes a fine logarithmic grid, then chooses the finite-$N$
threshold for that grid. Choosing $\tau$ inside the strict gap
$0.49254-I_{\rm Liu}$ gives $S_{\rm rec}(N)\le0.49254/\log N$
eventually, with the exact finite pair carrier unchanged. These are the
actual producers in
\sourcefile{MathlibNt/SieveTheory/Liu/PrimePairs/LiuPrimePairTransfer.lean}{prime-pair transfer}.

To pass to the genuine main mass, fix $\kappa\ge0$ and put
\[\begin{aligned}
 \Delta_{\kappa}(x)&=\operatorname{li}_{\kappa}(x)-\frac{x}{\log x},\\
 H_{\kappa}(N)&=\sum_a b(a)\operatorname{li}_{\kappa}(N/a)
       =N S_{\rm rec}(N)+\sum_{(r,s)\in\mathcal L}\Delta_{\kappa}(N/(rs)).
\end{aligned}\]
The equality is exact reindexing by the unique prime pair. Integration by
parts gives $|\Delta_{\kappa}(x)|\le C_{\Delta,\kappa}x/(\log x)^2$ eventually.
For every source pair, $N/(rs)\ge N^{1/3}$, hence
$\log(N/(rs))\ge(\log N)/3$. Mertens bounds the reciprocal pair mass by
a fixed constant $B_{\rm pair}$, so uniformly over the pair set,
\[
 \left|\sum_{(r,s)\in\mathcal L}\Delta_{\kappa}(N/(rs))\right|
 \le\frac{9C_{\Delta,\kappa}B_{\rm pair}N}{(\log N)^2}
 =o\!\left(\frac{N}{\log N}\right).
\]
Thus for every fixed $\kappa\ge0$ and $\eta>0$, eventually
\[
 0\le H_{\kappa}(N)\le(0.49254+\eta)\frac{N}{\log N}.
\]
The exact decomposition and summed correction are proved in
\sourcefile{MathlibNt/SieveTheory/Liu/Weights/LiuWeightMainSum.lean}{genuine-li correction};
\sourcefile{MathlibNt/SieveTheory/Selberg/Liu/LiuSelbergMainTermInstantiation.lean}{unconditional genuine-li mass}
combines this correction with the proved reciprocal-log bound.
The actual even-filter optimizer consumes that result in
\sourcefile{MathlibNt/SieveTheory/Selberg/Liu/LiuSelbergEvenAssembly.lean}{optimized main-term assembly}.
It fixes $\delta=\eta=10^{-7}$ and $0<\epsilon\le\epsilon_0=10^{-10}$;
the coefficient product satisfies
$(8+\delta)(0.49254+\eta)\le3.94033$.
Consequently the displayed main term $M_1=H_2(N)/\mathcal G$ satisfies
$M_1\le3.94033\mathcal X_N$ with a genuine, fixed coefficient margin.
\inputleannode{bp:chen-square-main-term}

The remaining distribution theorem concerns the convolution $b(a)$ with
primes.  For $A_1=\lfloor(\log N)^{2B}\rfloor+1$ and
$A_2=\lfloor N^{2/3}\rfloor$, define
\[
 E_q(\kappa)=\max_{l\in(\mathbb Z/q\mathbb Z)^\times}
 \left|\sum_{\substack{A_1<a\le A_2\\(a,q)=1}}b(a)
  \left(\pi(N;a,q,l)-\frac{\operatorname{li}_{\kappa}(N/a)}{\varphi(q)}\right)\right|.
\]
The absolute value is outside the $a$-sum.  Modulus one uses residue zero; modulus zero contributes zero.  At the fixed normalization $\kappa=2/\log2$,
the proved source theorem says that every $\sigma>0$ admits $C,B,N_0$ such that
$\sum_{q<N^{1/2}/(\log N)^B}E_q(\kappa)\le CN/(\log N)^\sigma$
for every $N\ge N_0$.  Its proof splits nonprincipal primitive conductors
into low and high ranges, uniformly in the changing source and cofactor,
and pays the principal character separately.  Two extra logarithms pay the
reciprocal-totient cofactor sum before this unweighted bound is obtained.
See the source-linked statement in result~\ref{bp:foundation-liupan}.
The Selberg expansion needs a weighted modulus sum.  Finite Cauchy gives
\[\begin{gathered}
 \left(\sum_{q<L}\mu(q)^2 3^{\omega(q)}E_q(\kappa)\right)^2\\
 \le C_9(\log(N+2))^9
       \left(C_\kappa N(1+\log N)^2\sum_{q<L}E_q(\kappa)\right),\\
 L=N^{1/2}/(\log N)^B.
\end{gathered}\]
Requesting $\sigma=2A+11$ therefore leaves any prescribed saving $A$.
\inputleannode{bp:chen-switched-weight-payment}
Source support identifies the interval sum with the full supported sum.
Replacing the strict source modulus range by the smaller closed range costs
$B\mapsto B+1$; changing the additive normalization to $\kappa=2$ is paid
separately.  This supplies the canonical coprime distribution bound.
In the square, $[d_1,d_2]\le R^2\le N^{1/2-\epsilon}$; regrouping coefficient
pairs gives the $3^{\omega(d)}$ majorant.  The noncoprime $a$-part has an
independent $O(N^{9/10}(\log N)^2)$ bound, so it too fits every fixed
inverse-log budget.  The resulting estimate is
$|\mathcal E_\lambda|\le C N/(\log N)^A$, uniformly for the eventual
admissible coefficient family, with $C$ fixed before $N$.
\inputleannode{bp:chen-square-remainder}

\subsection{Closing one budget on the original representation count}
Combine the square bound with the two finite corrections to obtain
\[
 T(N)\le3.94033\mathcal X_N+C\frac{N}{(\log N)^3}
\]
eventually for even $N$.  In the actual assembly $\epsilon=\epsilon_0$ and
$A=3$ are fixed first; the resulting constant is the sum of the square,
endpoint, and paper-modulus residual constants.
\inputleannode{bp:chen-triple-upper}
Now choose the weighted-lower margin $\eta=1/10000$.
The uniform bound $\mathfrak S_{\mathrm{Liu}}(N)\ge U>0$ makes the displayed remainder
at most $\mathcal X_N/10000$ eventually.  Intersect these two events with the
triple-bound event and $N\ge9$.  The finite bridge then gives
\[
 |G(N)|\ge W(N)-\tfrac12T(N)
 \ge\left(2.6408-10^{-4}-\tfrac12(3.94033+10^{-4})\right)\mathcal X_N
 =0.670485\mathcal X_N\ge0.67\mathcal X_N.
\]
This is the public count introduced at the start, with the same prime, square, and unit conventions.
\inputleannode{bp:chen-representation-lower-bound}
Since $\mathcal X_N>0$ for $N>1$, a positive finite cardinality supplies a prime
$p\in G(N)$ and hence $N=p+(N-p)$.  The implemented existence route uses
the same finite positivity lemma with the more generous margins $1/10$,
then extracts an eventual threshold $N_0$ and expands the almost-prime
predicate.  Its conclusion permits either a prime partner or $q=rs$ with
$r,s$ prime, including $r=s$.
\inputleannode{bp:chen-theorem}
